\documentclass[11pt]{article}
\usepackage[a4paper,margin=23mm]{geometry}
\usepackage[T1]{fontenc}
\usepackage{lmodern,amsmath,amssymb,amsthm,mathtools,microtype,booktabs}
\usepackage[colorlinks=true,linkcolor=blue,citecolor=blue,urlcolor=blue]{hyperref}
\newcommand{\EE}{\mathbb E}
\newcommand{\tr}{\operatorname{tr}}
\newcommand{\PP}{\mathcal P_h}
\newcommand{\FF}{\mathcal F}
\newcommand{\BB}{\mathcal B}
\newcommand{\HH}{H^{(q,C)}}
\newtheorem{theorem}{Theorem}
\newtheorem{lemma}{Lemma}
\title{Reconstructing the full bilocal fluctuation\\
\large A strong Volterra limit and explicit spectral amplitudes}
\author{Alan Day \and Chengheng\thanks{AI research assistant (OpenAI).}}
\date{7 November 2026\quad Local research draft}
\begin{document}\maketitle
\begin{abstract}
The previously constructed joint finite-time memory kernel is now
connected to the original singlet bilocal. A Hermitian dilation of
the energy-pinched memory makes its background propagator contractive,
including degenerate energy blocks. Interpolation between the collective
second-moment bound and the existing force eighth moment controls the
return terms, giving an $O(p^{-1/3})$ strong reconstruction error on
fixed time intervals. The charged spectral amplitude simplifies to a
single centered Hermite coefficient; the neutral amplitude is a regular
second divided difference of the exponential. Their common positive
spectral matrix gives the complete bilocal covariance. This is a result
for the specified microscopic model; the bulk mass and position map
remains to be constructed.
\end{abstract}

\section{Model and the actual singlet observable}
Use standard CAR and the overlapping complex Gaussian $p$-body
Hamiltonian, with the same original disorder across charge sectors.
The state is normalized annealed full-Fock trace at $\beta=\mu=0$.
Take even $N$, $p\to\infty$, $\lambda_N=p^2/N\to\lambda>0$, fixed
$J_s$, $\hbar=1$, and string time. Write
\begin{equation}
 C=J_s^2/\lambda,\quad q=e^{-4\lambda},\quad
 b_{N,m}=2^{-N}\binom Nm.
\end{equation}
In a sector $m>0$, put $h=H_m$, $K=I_N\otimes H_{m-1}$ and
\begin{equation}
 V\psi=m^{-1/2}\sum_i|i\rangle c_i\psi,\quad
 P=VV^\dagger,\quad R=(1-P)KV,\quad D=(1-P)K(1-P).
\end{equation}
On $P\oplus(1-P)$, $K$ has system block $\alpha h$, where
$\alpha=1-p/m$. Its compressed propagator and memory are
\begin{equation}
 F(t)=V^\dagger e^{iKt}V,\quad M(t)=R^\dagger e^{iDt}R,\quad
 F'=i\alpha hF-M*F,\quad F(0)=1.
\end{equation}
Convolution is ordered: $(A*B)(t)=\int_0^t A(t-s)B(s)ds$.
The center-time-zero bilocal and its dynamic part are
\begin{equation}
 S_c(t)=m e^{-iht/2}F(t)e^{-iht/2},\quad
 \BB_N(t)=(1-\PP)S_c(t),\quad \PP Y=\sum_eP_eYP_e.
\end{equation}
The sum keeps whole energy eigenspaces. Inactive $m<p$ sectors
have zero dynamic field, including $m=0$. We reconstruct this
unnormalized fluctuation, whose limit is nontrivial even though
$S_c/m\to1$ \cite{Channel}.

\newpage
\section{A contractive background and an exact difference}
Set $M_0=\PP M$ and $\delta M=M-M_0$. Let $F_0$ solve the same
Volterra equation as $F$, with $M_0$ in place of $M$.
It is not enough to assume that a pinched memory has a stable
propagator; here this follows from an explicit dilation.

For every energy projector $P_e$, introduce a separate copy of
the original bath and define
\begin{equation}
 R_0\psi=(RP_e\psi)_e,\qquad
 D_0=\bigoplus_e D,\qquad
 K_0=\begin{pmatrix}\alpha h&R_0^\dagger\\R_0&D_0\end{pmatrix}.
\end{equation}
Then $R_0^\dagger e^{iD_0t}R_0=M_0(t)$. Consequently $F_0$ is the
system compression of $e^{iK_0t}$ and
\begin{equation}
 \|F(t)\|_\infty\le1,\qquad \|F_0(t)\|_\infty\le1.
 \label{eq:contract}
\end{equation}
Each system energy block interacts only with its own bath copy.
Thus $F_0$ commutes with $h$, but need not be scalar inside a
degenerate block. Kernels $M_0(s)$ and $M_0(t)$ need not commute.

For $\operatorname{Re}z>0$, the compressed resolvents are
\begin{equation}
 \widehat F=(z-i\alpha h+\widehat M)^{-1},\qquad
 \widehat F_0=(z-i\alpha h+\widehat M_0)^{-1}.
\end{equation}
The resolvent identity gives the exact ordered relations
\begin{align}
 F-F_0&=-F_0*\delta M*F,\label{eq:exact}\\
 &=-F_0*\delta M*F_0+F_0*\delta M*F_0*\delta M*F.
\end{align}
No infinite perturbation expansion is used. The preceding memory
field is
\begin{equation}
 \FF_N(s)=(1-\PP)m e^{-ihs/2}M(s)e^{-ihs/2},\qquad
 Q_N(s):=m\delta M(s)=e^{ihs/2}\FF_N(s)e^{ihs/2}.
 \label{eq:Q}
\end{equation}
The already proved uniform bound is
\begin{equation}
 \sup_N\sup_{|s|\le S}\|Q_N(s)\|_{2,\mathrm{ann}}<\infty
 \quad(S<\infty).\label{eq:collective}
\end{equation}
In particular \eqref{eq:exact} identifies the natural scale without
discarding the finite return term. It also implies
$\|m(F-F_0)(t)\|_2\le t^2\sup_{s\le t}\|Q_N(s)\|_2/2$.

\section{Which norm is being used}
For sector operators let
\begin{equation}
 \|X\|_{r,A}^r=\sum_{m\in A}b_{N,m}\EE\tr_m|X_m|^r.
\end{equation}
The Gaussian expectation and charge sum form one finite trace;
Schatten H\"older and interpolation apply in this algebra.
Rectangular factors are normalized by the same source dimension.
All estimates below are uniform on a fixed real-time interval;
constants can depend on that interval and a positive compact
interval containing $\lambda_N$.

\newpage
\section{Interpolation controls the collective return}
The existing force square-function estimate \cite{Return} gives,
for $G=R^\dagger R$ and $\mathcal A_m=\sum_i[H,c_i]^\dagger[H,c_i]|_m$,
\begin{align}
 0\le G&\le\mathcal A_m/m,\\
 \EE\tr_m\mathcal A_m^k&\le(2k-1)!!(Nv_N)^k,\qquad Nv_N=O(p).
 \label{eq:force}
\end{align}
Trace eigenvalue monotonicity, not operator monotonicity of the
power map, gives $\EE\tr_mG^k\le m^{-k}\EE\tr_m\mathcal A_m^k$.
Rectangular H\"older yields
\begin{equation}
 \|M(s)\|_{k,m,\mathrm{ann}}\le\|G\|_{k,m,\mathrm{ann}}
 \le ((2k-1)!!)^{1/k}Nv_N/m.
\end{equation}
Pinching contracts every Schatten norm, so $M_0$ has the same bound
and $\delta M$ at most twice this bound.

Take the good charge set $A_N=\{m\ge N/4\}$. On it, the stronger
second moment follows from \eqref{eq:collective}, whereas the
eighth moment follows from \eqref{eq:force}:
\begin{equation}
 \|\delta M(s)\|_{2,A_N}=O(p^{-2}),\qquad
 \|\delta M(s)\|_{8,A_N}=O(p^{-1}).
\end{equation}
Thus the standard $2$--$8$ interpolation inequality gives
\begin{equation}
 \|\delta M(s)\|_{4,A_N}
 \le\|\delta M(s)\|_{2,A_N}^{1/3}
       \|\delta M(s)\|_{8,A_N}^{2/3}
 =O(p^{-4/3}).\label{eq:interp}
\end{equation}
This does not assert a uniform fourth moment for $m\delta M$.

Let $U(t)=e^{iht}$ and $U_\alpha(t)=e^{i\alpha ht}$.
Using \eqref{eq:contract}, the exact integral equations
$F=U_\alpha-U_\alpha*M*F$ and
$F_0=U_\alpha-U_\alpha*M_0*F_0$ imply
\begin{equation}
 \|F-U\|_{4,A_N}+\|F_0-U\|_{4,A_N}=O_S(p^{-1}).\label{eq:free}
\end{equation}
Indeed $M,M_0$ have fourth norm $O(p^{-1})$, while
$\|(\alpha-1)h\|_{4,A_N}\le(4p/N)\|h\|_4=O(p^{-1})$.
The uniform fourth energy moment is already available in this model.

In the exact first line of \eqref{eq:exact}, the integrand error
after replacing both propagators by $U$ is
\begin{equation}
 (F_0-U)\,m\delta M\,F+U\,m\delta M\,(F-U).
\end{equation}
H\"older, $m\le N=O(p^2)$ and \eqref{eq:interp}--\eqref{eq:free}
bound its second norm on $A_N$ by $O(p^{-1/3})$.
On the remaining active set $p\le m<N/4$,
\begin{equation}
 \sup_{|s|\le S}\|Q_N(s)\|_{2,A_N^c}
 \le Kp\,e^{-N/16},
\end{equation}
using the sector bound $\|m\delta M\|_2\le2\sqrt{3}\,Nv_N$
and $\sum_{m<N/4}b_{N,m}\le e^{-N/8}$.
There the propagator differences have operator norm at most two.
This treats all charges without a new four-source moment theorem.

\newpage
\section{Strong reconstruction of the entire bilocal}
Define, for $t\ge0$,
\begin{equation}
 (\mathcal V_h\FF_N)(t)=
 -\int_0^t ds\int_{-(t-s)/2}^{(t-s)/2}du\,
 e^{ihu}\FF_N(s)e^{-ihu}.\label{eq:filter}
\end{equation}
For negative times use the adjoint relation. This is the whole
memory field, with its charged, neutral and mixed components retained.

\begin{theorem}
For each fixed $S<\infty$, in the stated model and limit,
\begin{equation}
 \sup_{|t|\le S}\|\BB_N(t)-(\mathcal V_h\FF_N)(t)\|_{2,\mathrm{ann}}
 \le K_Sp^{-1/3}+K_Sp e^{-N/16}.\label{eq:strong}
\end{equation}
Both fields have uniformly bounded second norms on this interval.
\end{theorem}
\begin{proof}
The background $m e^{-iht/2}F_0(t)e^{-iht/2}$ is removed exactly
by $1-\PP$. Apply the preceding integrand bounds to
\eqref{eq:exact}, whose two-time simplex has area $t^2/2$.
After replacing its propagators by $U$, put their two times equal
to $a,b\ge0$, with $a+b=t-s$. Equation \eqref{eq:Q} combines all
phases into $e^{ihu}\FF_N(s)e^{-ihu}$, where $u=(a-b)/2$ and
$da=du$. This gives \eqref{eq:filter} with no extra factor.
Its pinching is zero already, since
$\PP(e^{ihu}\FF_N(s)e^{-ihu})=0$.
Conjugation preserves the second norm, so
$\|\mathcal V_h\FF_N(t)\|_2\le t^2\sup_{s\le t}\|\FF_N(s)\|_2/2$.
The uniform bound and \eqref{eq:strong} complete the proof.
\end{proof}

Consequently the two covariance functions differ by $O_S(p^{-1/3})$
uniformly in all external real center times $T$. This is a comparison
of two finite-$N$ covariances. Their limit below is asserted for
fixed $T$; no growing-time spectral limit is taken.

Let $\mathcal K_\FF(T;s,r)$ be the complete memory covariance from
\cite{Joint}. Dominated convergence and \eqref{eq:strong} give
\begin{align}
 \mathcal K_\BB(T;t,v)
 ={}&\int_0^t ds\int_0^v dr
 \int_{-(t-s)/2}^{(t-s)/2}du
 \int_{-(v-r)/2}^{(v-r)/2}dw\,
 \mathcal K_\FF(T+w-u;s,r).\label{eq:cov}
\end{align}
Here $t,v\ge0$, with other signs determined by adjoints.
The old finite-time memory construction has thus been connected
to the actual unnormalized singlet fluctuation, rather than only
to its first few Taylor coefficients.

\newpage
\section{Explicit charged and neutral spectral amplitudes}
Recall the common positive two-energy matrix $\nu_{ab}$ for
$\mathcal X_a\in\{J_d\ (d\ge2),N_j\ (j\ge1)\}$ from \cite{Joint}.
Here $N_j$ denotes a neutral residual field, not fermion number.
Energies $E,E'$ label input/column and output/row. Put
\begin{equation}
 e=(E+E')/2,\quad\omega=(E'-E)/2,\quad
 \chi_L(E'-E)=\frac{\sin(\omega L)}\omega,\quad\chi_L(0)=L.
\end{equation}
With $P_d=[d]_q!$ the monic Hermite coefficients are
\begin{equation}
 c_d(s)=\frac{\int e^{isx}\HH_d(x)d\rho_q(x)}{C^dP_d}.
\end{equation}
The centered memory amplitudes are
$e^{-ies}D_d(s;E,E')$ and $2Ce^{-ies}c_j(s)$, where
$D_d=-c_d''+2ie c_d'+EE'c_d$. Therefore \eqref{eq:filter} gives
\begin{equation}
 b_a(t;E,E')=-\int_0^t\chi_{t-s}(E'-E)e^{-ies}f_a(s;E,E')ds.
 \label{eq:bint}
\end{equation}

For the charged part let $g_d(s)=e^{-ies}c_d(s)$. Direct differentiation
gives $e^{-ies}D_d=-g_d''-\omega^2g_d$. The Green-function identity is
\begin{equation}
 -\int_0^t\chi_{t-s}e^{-ies}D_d(s)ds
 =g_d(t)-\cos(\omega t)g_d(0)-\chi_tg_d'(0).
\end{equation}
Since $c_d(0)=c_d'(0)=0$ for $d\ge2$, the charged amplitude is simply
\begin{equation}
 \boxed{b_{J_d}(t;E,E')=e^{-iet}c_d(t).}\label{eq:bcharged}
\end{equation}
The zero initial data are essential; this identity is not asserted
for the discarded $d=0,1$ channels.

For the neutral part use the standard second divided difference
of $f_t(x)=e^{itx}$ \cite{DLMF}. Integrating the scalar Green function,
or checking its initial-value equation, gives
\begin{equation}
 \boxed{b_{N_j}(t;E,E')=
 \frac{2Ce^{-iet}}{C^jP_j}\int\HH_j(x)
 [E,E',x]f_t\,d\rho_q(x).}\label{eq:bneutral}
\end{equation}
The apparent denominators in the distinct-node interpolation formula
are removable. A definition valid also for coincident real nodes is
\begin{equation}
 [x_0,x_1,x_2]f_t=(it)^2
 \int_{a,b\ge0,\ a+b\le1}
 e^{it[ax_0+bx_1+(1-a-b)x_2]}da\,db.
\end{equation}
Thus $|[x_0,x_1,x_2]f_t|\le t^2/2$, and at three coincident nodes
it equals $-t^2e^{itx_0}/2$. Two applications of the fundamental
theorem of calculus establish the integral identity, so no complex
mean-value point is assumed.

\newpage
\section{The resulting full covariance and its first coefficients}
The complete bilocal answer is
\begin{equation}
 \boxed{\mathcal K_\BB(T;t,v)=
 \sum_{a,b}\int e^{iT(E'-E)}
 \overline{b_a(t;E,E')}b_b(v;E,E')d\nu_{ab}(E,E').}
 \label{eq:answer}
\end{equation}
All charged/neutral mixed entries of the same original-disorder
matrix are included. The time and field sums can be integrated
termwise: on $0\le t\le S$, $|\chi_{t-s}|\le t-s$, and the previously
proved memory tails are absolutely summable in the field Hilbert
norm and in spectral total variation. The Volterra map has norm
at most $S^2/2$ on these time families. This transfers those tails
to \eqref{eq:answer} without an unproved uniform fourth moment.

Each finite spectral block has an $L^2(\rho_q\otimes\rho_q)$ density.
The full time-dependent spectral measure is absolutely continuous
by the total-variation limit, and has no diagonal static mass.
Only $L^1$ regularity is claimed for the infinite density. Positivity
of $\nu$ yields a nonnegative equal-time norm and
$\mathcal K_\BB(T;t,v)^*=\mathcal K_\BB(-T;v,t)$.

Writing $\FF(s)=A_0+is(A_1+B_1)+O(s^2)$ gives
\begin{align}
 \BB(t)&=-\frac{t^2}{2}J_2
 -\frac{it^3}{6}\left[J_3-\frac32\{H,J_2\}+2CN_1\right]
 +O_{L^2}(t^4),\\
 \mathcal K_\BB(T;t,v)&=\frac{t^2v^2}{4}\mathcal K_{J_2}(T)
 +O\bigl(|t|^3|v|^2+|t|^2|v|^3\bigr).
\end{align}
These are expansions of the established limiting field. They recover
the earlier $-t^2Y_-/2$ bilocal curvature. The nonzero memory channel
$\nu_{A_0B_2}$ is retained in \eqref{eq:answer}; no independence
assumption sets its contribution to zero.

\section{What physical information is still needed}
This stage determines the dynamic bilocal covariance in the stated
infinite-temperature double-scaled model. Its Gram metric can now be
tested against candidate bulk constructions. That metric by itself
does not supply canonical bosonic dressing or a light-front Hamiltonian.
The existing Itakura construction \cite{Itakura} fixes a vacuum,
normal ordering, polarization and physical boson subspace before its
meson projection. Our finite normalized trace obeys
$\tr[A,B]=0$ exactly, so a nonzero vacuum CCR expectation cannot be
read off in the same state with convergent commutator moments.

The next task is to identify the relevant algebra and state, and a
microscopic spatial/internal-momentum generator. The energy difference
$E'-E$, the frequency conjugate to memory time, and a meson's internal
quark momentum fraction are presently different variables. A genuine
$P^-$ and its interactions must be derived before comparing with the
already known bulk mass operator. The present result does not assign
particle masses to the spectral frequencies or establish finite-$G$
gravitational four-point dynamics.

\newpage
\section{Independent checks and reproducibility}
Two new scientific outputs are exclusive-create and retained unchanged.
The final evidence binds their input hashes, parameters and dependency
versions, predecessor files, both compilation passes and every final
rendered page. No earlier scientific main is rerun.
\begin{itemize}
\item Finite Volterra checks: two nontrivial CAR sectors
$(N,p,m)=(4,2,3),(5,2,3)$, one zero-memory $m=p$ control, and
an independent system with a two-dimensional degenerate energy block.
There are eight Laplace-resolvent cases, sixteen time-domain integral
checks, twenty-four contraction checks and twelve curvature values.
The maximum second-order resolvent error is $2.69\times10^{-15}$.
Time quadratures use orders 8 and 14. Nonzero finite return terms
are saved; these small systems are not a double-scaling fit.
\item The degenerate example has noncommuting $M_0(0.3),M_0(0.8)$
(normalized Hilbert--Schmidt commutator about $0.0331$).
Replacing whole eigenspaces by separate eigenvectors changes the
background memory, by about $0.259$ in the recorded test.
This checks that block degeneracy is treated as part of the definition.
\item Spectral-filter checks: 1,260 charged and 1,440 neutral exact
rational coefficient identities through degree 14. Two physical
$\lambda$ values, four signed times, ordinary/repeated/near-repeated
energy pairs and six field degrees give 144 floating field cases.
Independent spectral/time sizes $(40,24)$ and $(64,36)$ agree to
$2.54\times10^{-15}$. The time-integral and explicit-amplitude
forms differ by at most $2.40\times10^{-15}$.
\item Twenty additional repeated-node cases compare a $3\times3$
triangular matrix exponential with a simplex integral. The triple-node
value and $t^2/2$ bound are checked separately. All floating comparisons
retain their errors and are not interval certificates.
\end{itemize}
The strong-return theorem uses the analytical norm estimates, not
extrapolation from these tests. This stage remains local, with no Git
publication or external issue reply.

\begin{thebibliography}{9}
\bibitem{Channel} A. Day and Chengheng,
\emph{The compressed evolution and memory of the complete bilocal},
local project stage \texttt{susskind-bilocal-channel} (2026).
\bibitem{Return} A. Day and Chengheng,
\emph{Two-return reduction of finite memory}, local project stage
\texttt{susskind-memory-return-expansion} (2026), force high moments;
and \texttt{susskind-memory-uniform-bound}, collective second norm.
\bibitem{Joint} A. Day and Chengheng,
\emph{The joint finite-time memory kernel}, local project stage
\texttt{susskind-memory-joint-regression} (2026).
\bibitem{DLMF} NIST Digital Library of Mathematical Functions,
\href{https://dlmf.nist.gov/3.3.iii}{Sec.~3.3(iii)},
Eqs.~3.3.34--3.3.37, divided differences and analytic continuation.
\bibitem{Itakura} K. Itakura,
\emph{Boson Expansion Methods in (1+1)-dimensional Light-Front QCD},
\href{https://arxiv.org/abs/hep-th/9604032}{arXiv:hep-th/9604032},
Secs.~4--5. The source record distinguishes its light-front conventions
and vacuum from the present model.
\end{thebibliography}
\end{document}
